\iftestbuild % TEST-ONLY-BEGIN
\setcounter{section}{4}\setcounter{theorem}{63}
\fi % TEST-ONLY-END
% Source: book pp. 270--279 (PDF 284--293); no solutions.
% \axeqno: see 7C.tex
\DeclareRobustCommand\axeqno{\refstepcounter{theorem}\tag*{\thetheorem}}
\section{Singular Value Decomposition}

\subsection{Singular Values}

We will need the following result in this section.

\begin{result}[Properties of $T^*T$]\label{ax:7.64}
Suppose $T\in\Lin(V,W)$. Then
\begin{parts}
\item $T^*T$ is a positive operator on $V$;
\item $\nul T^*T=\nul T$;
\item $\range T^*T=\range T^*$;
\item $\dim\range T=\dim\range T^*=\dim\range T^*T$.
\end{parts}
\end{result}

\begin{proof}\leavevmode
\begin{parts}
\item We have
  \[ (T^*T)^* = T^*(T^*)^* = T^*T. \]
  Thus $T^*T$ is self-adjoint.

  If $v\in V$, then
  \[ \bigl\langle (T^*T)v, v\bigr\rangle = \ip{T^*(Tv)}{v} = \ip{Tv}{Tv}
     = \norm{Tv}^2 \ge 0. \]
  Thus $T^*T$ is a positive operator.
\item First suppose $v\in\nul T^*T$. Then
  \[ \norm{Tv}^2 = \ip{Tv}{Tv} = \ip{T^*Tv}{v} = \ip{0}{v} = 0. \]
  Thus $Tv=0$, proving that $\nul T^*T\subseteq\nul T$.

  The inclusion in the other direction is clear, because if $v\in V$ and
  $Tv=0$, then $T^*Tv=0$.

  Thus $\nul T^*T=\nul T$, completing the proof of (b).
\item We already know from (a) that $T^*T$ is self-adjoint. Thus
  \[ \range T^*T = (\nul T^*T)^\perp = (\nul T)^\perp = \range T^*, \]
  where the first and last equalities come from \ref{ax:7.6} and the second
  equality comes from (b).
\item To verify the first equation in (d), note that
  \[ \dim\range T = \dim(\nul T^*)^\perp = \dim W-\dim\nul T^* = \dim\range T^*, \]
  where the first equality comes from \ref{ax:7.6}(d), the second equality
  comes from \ref{ax:6.51}, and the last equality comes from the fundamental
  theorem of linear maps (\ref{ax:3.21}).

  The equality $\dim\range T^*=\dim\range T^*T$ follows from (c).\qedhere
\end{parts}
\end{proof}

The eigenvalues of an operator tell us something about the behavior of the
operator. Another collection of numbers, called the singular values, is also
useful. Eigenspaces and the notation $E$ (used in the examples) were defined in
\ref{ax:5.52}.

\begin{definition}[Singular values]\label{ax:7.65}
Suppose $T\in\Lin(V,W)$. The \emph{singular values} of $T$ are the nonnegative
square roots of the eigenvalues of $T^*T$, listed in decreasing order, each
included as many times as the dimension of the corresponding eigenspace of
$T^*T$.
\end{definition}

\begin{example}[Singular values of an operator on $\F^4$]\label{ax:7.66}
Define $T\in\Lin(\F^4)$ by $T(z_1,z_2,z_3,z_4)=\allowbreak(0,3z_1,2z_2,-3z_4)$. A
calculation shows that
\[ T^*T(z_1,z_2,z_3,z_4) = (9z_1,4z_2,0,9z_4), \]
as you should verify. Thus the standard basis of $\F^4$ diagonalizes $T^*T$,
and we see that the eigenvalues of $T^*T$ are $9$, $4$, and $0$. Also, the
dimensions of the eigenspaces corresponding to the eigenvalues are
\[ \dim E(9,T^*T)=2 \quad\text{and}\quad \dim E(4,T^*T)=1 \quad\text{and}\quad
   \dim E(0,T^*T)=1. \]
Taking nonnegative square roots of these eigenvalues of $T^*T$ and using
dimension information from above, we conclude that the singular values of $T$
are $3,3,2,0$.

The only eigenvalues of $T$ are $-3$ and $0$. Thus in this case, the collection
of eigenvalues did not pick up the number $2$ that appears in the definition
(and hence the behavior) of $T$, but the list of singular values does include
$2$.
\end{example}

\begin{example}[Singular values of a linear map from $\F^4$ to
$\F^3$]\label{ax:7.67}
Suppose $T\in\Lin(\F^4,\F^3)$ has matrix (with respect to the standard bases)
\[ \begin{pmatrix} 0&0&0&-5\\ 0&0&0&0\\ 1&1&0&0 \end{pmatrix}. \]
You can verify that the matrix of $T^*T$ is
\[ \begin{pmatrix} 1&1&0&0\\ 1&1&0&0\\ 0&0&0&0\\ 0&0&0&25 \end{pmatrix} \]
and that the eigenvalues of the operator $T^*T$ are $25,2,0$, with
$\dim E(25,T^*T)=1$, $\dim E(2,T^*T)=1$, and $\dim E(0,T^*T)=2$. Thus the
singular values of $T$ are $5,\sqrt2,0,0$.
\end{example}

See Exercise~2 for a characterization of the positive singular values.

\begin{result}[Role of positive singular values]\label{ax:7.68}
Suppose that $T\in\Lin(V,W)$. Then
\begin{parts}
\item $T$ is injective $\iff$ $0$ is not a singular value of $T$;
\item the number of positive singular values of $T$ equals $\dim\range T$;
\item $T$ is surjective $\iff$ number of positive singular values of $T$
  equals $\dim W$.
\end{parts}
\end{result}

\begin{proof}
The linear map $T$ is injective if and only if $\nul T=\{0\}$, which happens if
and only if $\nul T^*T=\{0\}$ [by \ref{ax:7.64}(b)], which happens if and only
if $0$ is not an eigenvalue of $T^*T$, which happens if and only if $0$ is not
a singular value of $T$, completing the proof of (a).

The spectral theorem applied to $T^*T$ shows that $\dim\range T^*T$ equals the
number of positive eigenvalues of $T^*T$ (counting repetitions). Thus
\ref{ax:7.64}(d) implies that $\dim\range T$ equals the number of positive
singular values of $T$, proving (b).

Use (b) and \ref{ax:2.39} to show that (c) holds.
\end{proof}

The table below compares eigenvalues with singular values.

\begin{center}\small
\begin{tabular}{|p{.44\linewidth}|p{.44\linewidth}|}
\hline
\multicolumn{1}{|c|}{\textbf{list of eigenvalues}} &
\multicolumn{1}{c|}{\textbf{list of singular values}}\\
\hline
context: vector spaces & context: inner product spaces\\
\hline
defined only for linear maps from a vector space to itself &
defined for linear maps from an inner product space to a possibly different
inner product space\\
\hline
can be arbitrary real numbers (if $\F=\R$) or complex numbers (if $\F=\C$) &
are nonnegative numbers\\
\hline
can be the empty list if $\F=\R$ & length of list equals dimension of domain\\
\hline
includes $0$ $\iff$ operator is not invertible &
includes $0$ $\iff$ linear map is not injective\\
\hline
no standard order, especially if $\F=\C$ & always listed in decreasing order\\
\hline
\end{tabular}
\end{center}

The next result nicely characterizes isometries in terms of singular values.

\begin{result}[Isometries characterized by having all singular values
equal~$1$]\label{ax:7.69}
Suppose that $S\in\Lin(V,W)$. Then
\[ S \text{ is an isometry} \iff \text{all singular values of $S$ equal $1$}. \]
\end{result}

\begin{proof}
We have
\begin{align*}
S \text{ is an isometry} &\iff S^*S=I\\
  &\iff \text{all eigenvalues of $S^*S$ equal $1$}\\
  &\iff \text{all singular values of $S$ equal $1$},
\end{align*}
where the first equivalence comes from \ref{ax:7.49} and the second equivalence
comes from the spectral theorem (\ref{ax:7.29} or \ref{ax:7.31}) applied to the
self-adjoint operator $S^*S$.
\end{proof}

\subsection{SVD for Linear Maps and for Matrices}

\marginnote[-12mm]{The singular value decomposition is useful in computational linear
algebra because good techniques exist for approximating eigenvalues and
eigenvectors of positive operators such as $T^*T$, whose eigenvalues and
eigenvectors lead to the singular value decomposition.}%
The next result shows that every linear map from $V$ to $W$ has a remarkably
clean description in terms of its singular values and orthonormal lists in $V$
and $W$. In the next section we will see several important applications of the
singular value decomposition (often called the SVD).

\begin{result}[Singular value decomposition]\label{ax:7.70}
Suppose $T\in\Lin(V,W)$ and the positive singular values of $T$ are
$s_1,\dots,s_m$. Then there exist orthonormal lists $e_1,\dots,e_m$ in $V$ and
$f_1,\dots,f_m$ in $W$ such that
\[ Tv = s_1\ip{v}{e_1}f_1+\cdots+s_m\ip{v}{e_m}f_m \axeqno\label{ax:7.71} \]
for every $v\in V$.
\end{result}

\begin{proof}
Let $s_1,\dots,s_n$ denote the singular values of $T$ (thus $n=\dim V$).
Because $T^*T$ is a positive operator [see \ref{ax:7.64}(a)], the spectral
theorem implies that there exists an orthonormal basis $e_1,\dots,e_n$ of $V$
with
\[ T^*Te_k = s_k^2e_k \axeqno\label{ax:7.72} \]
for each $k=1,\dots,n$.

For each $k=1,\dots,m$, let
\[ f_k = \frac{Te_k}{s_k}. \axeqno\label{ax:7.73} \]
If $j,k\in\{1,\dots,m\}$, then
\[ \ip{f_j}{f_k} = \frac{1}{s_js_k}\ip{Te_j}{Te_k}
   = \frac{1}{s_js_k}\ip{e_j}{T^*Te_k} = \frac{s_k}{s_j}\ip{e_j}{e_k}
   = \begin{cases} 0 & \text{if } j\ne k,\\ 1 & \text{if } j=k. \end{cases} \]
Thus $f_1,\dots,f_m$ is an orthonormal list in $W$.

If $k\in\{1,\dots,n\}$ and $k>m$, then $s_k=0$ and hence $T^*Te_k=0$ (by
\ref{ax:7.72}), which implies that $Te_k=0$ [by \ref{ax:7.64}(b)].

Suppose $v\in V$. Then
\begin{align*}
Tv &= T(\ip{v}{e_1}e_1+\cdots+\ip{v}{e_n}e_n)\\
   &= \ip{v}{e_1}Te_1+\cdots+\ip{v}{e_m}Te_m\\
   &= s_1\ip{v}{e_1}f_1+\cdots+s_m\ip{v}{e_m}f_m,
\end{align*}
where the last index in the first line switched from $n$ to $m$ in the second
line because $Te_k=0$ if $k>m$ (as noted in the paragraph above) and the third
line follows from \ref{ax:7.73}. The equation above is our desired result.
\end{proof}

Suppose $T\in\Lin(V,W)$, the positive singular values of $T$ are
$s_1,\dots,s_m$, and $e_1,\dots,e_m$ and $f_1,\dots,f_m$ are as in the singular
value decomposition \ref{ax:7.70}. The orthonormal list $e_1,\dots,e_m$ can be
extended to an orthonormal basis $e_1,\dots,e_{\dim V}$ of $V$ and the
orthonormal list $f_1,\dots,f_m$ can be extended to an orthonormal basis
$f_1,\dots,f_{\dim W}$ of $W$. The formula \ref{ax:7.71} shows that
\[ Te_k = \begin{cases} s_kf_k & \text{if } 1\le k\le m,\\
   0 & \text{if } m<k\le\dim V. \end{cases} \]
Thus the matrix of $T$ with respect to the orthonormal bases
$(e_1,\dots,e_{\dim V})$ and $(f_1,\dots,f_{\dim W})$ has the simple form
\[ \Mat\bigl(T,(e_1,\dots,e_{\dim V}),(f_1,\dots,f_{\dim W})\bigr)_{j,k}
   = \begin{cases} s_k & \text{if } 1\le j=k\le m,\\
   0 & \text{otherwise}. \end{cases} \]

If $\dim V=\dim W$ (as happens, for example, if $W=V$), then the matrix
described in the paragraph above is a diagonal matrix. If we extend the
definition of diagonal matrix as follows to apply to matrices that are not
necessarily square, then we have proved the wonderful result that every linear
map from $V$ to $W$ has a diagonal matrix with respect to appropriate
orthonormal bases.

\begin{definition}[Diagonal matrix]\label{ax:7.74}
An $M$-by-$N$ matrix $A$ is called a \emph{diagonal matrix} if all entries of
the matrix are $0$ except possibly $A_{k,k}$ for $k=1,\dots,\min(M,N)$.
\end{definition}

The table below compares the spectral theorem (\ref{ax:7.29} and
\ref{ax:7.31}) with the singular value decomposition (\ref{ax:7.70}).

\begin{center}\small
\begin{tabular}{|p{.44\linewidth}|p{.44\linewidth}|}
\hline
\multicolumn{1}{|c|}{\textbf{spectral theorem}} &
\multicolumn{1}{c|}{\textbf{singular value decomposition}}\\
\hline
describes only self-adjoint operators (when $\F=\R$) or normal operators (when
$\F=\C$) &
describes arbitrary linear maps from an inner product space to a possibly
different inner product space\\
\hline
produces a single orthonormal basis &
produces two orthonormal lists, one for domain space and one for range space,
that are not necessarily the same even when range space equals domain space\\
\hline
different proofs depending on whether $\F=\R$ or $\F=\C$ &
same proof works regardless of whether $\F=\R$ or $\F=\C$\\
\hline
\end{tabular}
\end{center}

The singular value decomposition gives us a new way to understand the adjoint
and the inverse of a linear map. Specifically, the next result shows that given
a singular value decomposition of a linear map $T\in\Lin(V,W)$, we can obtain
the adjoint of $T$ simply by interchanging the roles of the $e$'s and the
$f$'s (see \ref{ax:7.77}). Similarly, we can obtain the pseudoinverse
$T^\dagger$ (see \ref{ax:6.68}) of $T$ by interchanging the roles of the $e$'s
and the $f$'s and replacing each positive singular value $s_k$ of $T$ with
$1/s_k$ (see \ref{ax:7.78}).

Recall that the pseudoinverse $T^\dagger$ in \ref{ax:7.78} below equals the
inverse $T^{-1}$ if $T$ is invertible [see \ref{ax:6.69}(a)].

\begin{result}[Singular value decomposition of adjoint and
pseudoinverse]\label{ax:7.75}
Suppose $T\in\Lin(V,W)$ and the positive singular values of $T$ are
$s_1,\dots,s_m$. Suppose $e_1,\dots,e_m$ and $f_1,\dots,f_m$ are orthonormal
lists in $V$ and $W$ such that
\[ Tv = s_1\ip{v}{e_1}f_1+\cdots+s_m\ip{v}{e_m}f_m \axeqno\label{ax:7.76} \]
for every $v\in V$. Then
\[ T^*w = s_1\ip{w}{f_1}e_1+\cdots+s_m\ip{w}{f_m}e_m \axeqno\label{ax:7.77} \]
and
\[ T^\dagger w = \frac{\ip{w}{f_1}}{s_1}e_1+\cdots+\frac{\ip{w}{f_m}}{s_m}e_m
   \axeqno\label{ax:7.78} \]
for every $w\in W$.
\end{result}

\begin{proof}
If $v\in V$ and $w\in W$ then
\begin{align*}
\ip{Tv}{w} &= \bigl\langle s_1\ip{v}{e_1}f_1+\cdots+s_m\ip{v}{e_m}f_m,
  w\bigr\rangle\\
&= s_1\ip{v}{e_1}\ip{f_1}{w}+\cdots+s_m\ip{v}{e_m}\ip{f_m}{w}\\
&= \bigl\langle v, s_1\ip{w}{f_1}e_1+\cdots+s_m\ip{w}{f_m}e_m\bigr\rangle.
\end{align*}
This implies that
\[ T^*w = s_1\ip{w}{f_1}e_1+\cdots+s_m\ip{w}{f_m}e_m, \]
proving \ref{ax:7.77}.

To prove \ref{ax:7.78}, suppose $w\in W$. Let
\[ v = \frac{\ip{w}{f_1}}{s_1}e_1+\cdots+\frac{\ip{w}{f_m}}{s_m}e_m. \]
Apply $T$ to both sides of the equation above, getting
\begin{align*}
Tv &= \frac{\ip{w}{f_1}}{s_1}Te_1+\cdots+\frac{\ip{w}{f_m}}{s_m}Te_m\\
   &= \ip{w}{f_1}f_1+\cdots+\ip{w}{f_m}f_m\\
   &= P_{\range T}w,
\end{align*}
where the second line holds because \ref{ax:7.76} implies that $Te_k=s_kf_k$
if $k=1,\dots,m$, and the last line above holds because \ref{ax:7.76} implies
that $f_1,\dots,f_m$ spans $\range T$ and thus is an orthonormal basis of
$\range T$ [and hence \ref{ax:6.57}(i) applies]. The equation above, the
observation that $v\in(\nul T)^\perp$ [see Exercise~8(b)], and the definition
of $T^\dagger w$ (see \ref{ax:6.68}) show that $v=T^\dagger w$, proving
\ref{ax:7.78}.
\end{proof}

\begin{example}[Finding a singular value decomposition]\label{ax:7.79}
Define $T\in\Lin(\F^4,\F^3)$ by $T(x_1,x_2,x_3,x_4)=(-5x_4,0,x_1+x_2)$. We want
to find a singular value decomposition of $T$. The matrix of $T$ (with respect
to the standard bases) is
\[ \begin{pmatrix} 0&0&0&-5\\ 0&0&0&0\\ 1&1&0&0 \end{pmatrix}. \]
Thus, as discussed in Example~\ref{ax:7.67}, the matrix of $T^*T$ is
\[ \begin{pmatrix} 1&1&0&0\\ 1&1&0&0\\ 0&0&0&0\\ 0&0&0&25 \end{pmatrix}, \]
and the positive eigenvalues of $T^*T$ are $25,2$, with
$\dim E(25,T^*T)=\allowbreak1$ and $\dim E(2,T^*T)=1$. Hence the positive singular values of $T$ are
$5,\sqrt2$.

Thus to find a singular value decomposition of $T$, we must find an
orthonormal list $e_1,e_2$ in $\F^4$ and an orthonormal list $f_1,f_2$ in
$\F^3$ such that
\[ Tv = 5\ip{v}{e_1}f_1+\sqrt2\ip{v}{e_2}f_2 \]
for all $v\in\F^4$.

An orthonormal basis of $E(25,T^*T)$ is the vector $(0,0,0,1)$; an orthonormal
basis of $E(2,T^*T)$ is the vector
$\bigl(1/\sqrt2,1/\sqrt2,0,0\bigr)$. Thus, following the proof
of \ref{ax:7.70}, we take
\[ e_1=(0,0,0,1) \quad\text{and}\quad
   e_2=\Bigl(1/\sqrt2,1/\sqrt2,0,0\Bigr) \]
and
\[ f_1=\frac{Te_1}{5}=(-1,0,0) \quad\text{and}\quad
   f_2=\frac{Te_2}{\sqrt2}=(0,0,1). \]
Then, as expected, we see that $e_1,e_2$ is an orthonormal list in $\F^4$ and
$f_1,f_2$ is an orthonormal list in $\F^3$ and
\[ Tv = 5\ip{v}{e_1}f_1+\sqrt2\ip{v}{e_2}f_2 \]
for all $v\in\F^4$. Thus we have found a singular value decomposition of $T$.
\end{example}

The next result translates the singular value decomposition from the context
of linear maps to the context of matrices. Specifically, the following result
gives a factorization of an arbitrary matrix as the product of three nice
matrices. The proof gives an explicit construction of these three matrices in
terms of the singular value decomposition.

In the next result, the phrase ``orthonormal columns'' should be interpreted to
mean that the columns are orthonormal with respect to the standard Euclidean
inner product.

\begin{result}[Matrix version of SVD]\label{ax:7.80}
Suppose $A$ is a $p$-by-$n$ matrix of rank $m\ge1$. Then there exist a
$p$-by-$m$ matrix $B$ with orthonormal columns, an $m$-by-$m$ diagonal matrix
$D$ with positive numbers on the diagonal, and an $n$-by-$m$ matrix $C$ with
orthonormal columns such that
\[ A = BDC^*. \]
\end{result}

\begin{proof}
Let $T\colon\F^n\to\F^p$ be the linear map whose matrix with respect to the
standard bases equals $A$. Then $\dim\range T=m$ (by \ref{ax:3.78}). Let
\[ Tv = s_1\ip{v}{e_1}f_1+\cdots+s_m\ip{v}{e_m}f_m \axeqno\label{ax:7.81} \]
be a singular value decomposition of $T$. Let
\begin{align*}
B &= \text{the $p$-by-$m$ matrix whose columns are $f_1,\dots,f_m$},\\
D &= \text{the $m$-by-$m$ diagonal matrix whose diagonal entries are
  $s_1,\dots,s_m$},\\
C &= \text{the $n$-by-$m$ matrix whose columns are $e_1,\dots,e_m$}.
\end{align*}

Let $u_1,\dots,u_m$ denote the standard basis of $\F^m$. If
$k\in\{1,\dots,m\}$ then
\[ (AC-BD)u_k = Ae_k-B(s_ku_k) = s_kf_k-s_kf_k = 0. \]
Thus $AC=BD$.

Multiply both sides of this last equation by $C^*$ (the conjugate transpose of
$C$) on the right to get
\[ ACC^* = BDC^*. \]

Note that the rows of $C^*$ are the complex conjugates of $e_1,\dots,e_m$. Thus
if $k\in\{1,\dots,m\}$, then the definition of matrix multiplication shows that
$C^*e_k=u_k$; hence $CC^*e_k=e_k$. Thus $ACC^*v=Av$ for all
$v\in\spn(e_1,\dots,e_m)$.

If $v\in\bigl(\spn(e_1,\dots,e_m)\bigr)^\perp$, then $Av=0$ (as follows from
\ref{ax:7.81}) and $C^*v=0$ (as follows from the definition of matrix
multiplication). Hence $ACC^*v=Av$ for all
$v\in\bigl(\spn(e_1,\dots,e_m)\bigr)^\perp$.

Because $ACC^*$ and $A$ agree on $\spn(e_1,\dots,e_m)$ and on
$\bigl(\spn(e_1,\dots,e_m)\bigr)^\perp$, we conclude that $ACC^*=A$. Thus the
displayed equation above becomes
\[ A = BDC^*, \]
as desired.
\end{proof}

Note that the matrix $A$ in the result above has $pn$ entries. In comparison,
the matrices $B$, $D$, and $C$ above have a total of
\[ m(p+m+n) \]
entries. Thus if $p$ and $n$ are large numbers and the rank $m$ is considerably
less than $p$ and $n$, then the number of entries that must be stored on a
computer to represent $A$ is considerably less than $pn$.

% ------------------------------------------------------------------ exercises
\exercisesheading{7E}

\begin{exercise}
Suppose $T\in\Lin(V,W)$. Show that $T=0$ if and only if all singular values of
$T$ are $0$.
\end{exercise}

\begin{exercise}
Suppose $T\in\Lin(V,W)$ and $s>0$. Prove that $s$ is a singular value of $T$ if
and only if there exist nonzero vectors $v\in V$ and $w\in W$ such that
\[ Tv=sw \quad\text{and}\quad T^*w=sv. \]
\exnote{The vectors $v,w$ satisfying both equations above are called a
\textbf{Schmidt pair}. Erhard Schmidt introduced the concept of singular values
in 1907.}
\end{exercise}

\begin{exercise}
Give an example of $T\in\Lin(\C^2)$ such that $0$ is the only eigenvalue of $T$
and the singular values of $T$ are $5,0$.
\end{exercise}

\begin{exercise}
Suppose that $T\in\Lin(V,W)$, $s_1$ is the largest singular value of $T$, and
$s_n$ is the smallest singular value of $T$. Prove that
\[ \{\norm{Tv} : v\in V \text{ and } \norm{v}=1\} = [s_n,s_1]. \]
\end{exercise}

\begin{exercise}
Suppose $T\in\Lin(\C^2)$ is defined by $T(x,y)=(-4y,x)$. Find the singular
values of $T$.
\end{exercise}

\begin{exercise}
Find the singular values of the differentiation operator
$D\in\Lin(\Poly_2(\R))$ defined by $Dp=p'$, where the inner product on
$\Poly_2(\R)$ is as in Example~\ref{ax:6.34}.
\end{exercise}

\begin{exercise}
Suppose that $T\in\Lin(V)$ is self-adjoint or that $\F=\C$ and $T\in\Lin(V)$ is
normal. Let $\lambda_1,\dots,\lambda_n$ be the eigenvalues of $T$, each
included in this list as many times as the dimension of the corresponding
eigenspace. Show that the singular values of $T$ are
$\abs{\lambda_1},\dots,\abs{\lambda_n}$, after these numbers have been sorted
into decreasing order.
\end{exercise}

\begin{exercise}
Suppose $T\in\Lin(V,W)$. Suppose $s_1\ge s_2\ge\cdots\ge s_m>0$ and
$e_1,\dots,e_m$ is an orthonormal list in $V$ and $f_1,\dots,f_m$ is an
orthonormal list in $W$ such that
\[ Tv = s_1\ip{v}{e_1}f_1+\cdots+s_m\ip{v}{e_m}f_m \]
for every $v\in V$.
\begin{parts}
\item Prove that $f_1,\dots,f_m$ is an orthonormal basis of $\range T$.
\item Prove that $e_1,\dots,e_m$ is an orthonormal basis of $(\nul T)^\perp$.
\item Prove that $s_1,\dots,s_m$ are the positive singular values of $T$.
\item Prove that if $k\in\{1,\dots,m\}$, then $e_k$ is an eigenvector of
  $T^*T$ with corresponding eigenvalue $s_k^2$.
\item Prove that
  \[ TT^*w = s_1^2\ip{w}{f_1}f_1+\cdots+s_m^2\ip{w}{f_m}f_m \]
  for all $w\in W$.
\end{parts}
\end{exercise}

\begin{exercise}
Suppose $T\in\Lin(V,W)$. Show that $T$ and $T^*$ have the same positive
singular values.
\end{exercise}

\begin{exercise}
Suppose $T\in\Lin(V,W)$ has singular values $s_1,\dots,s_n$. Prove that if $T$
is an invertible linear map, then $T^{-1}$ has singular values
\[ \frac{1}{s_n},\dots,\frac{1}{s_1}. \]
\end{exercise}

\begin{exercise}
Suppose that $T\in\Lin(V,W)$ and $v_1,\dots,v_n$ is an orthonormal basis of
$V$. Let $s_1,\dots,s_n$ denote the singular values of $T$.
\begin{parts}
\item Prove that $\norm{Tv_1}^2+\cdots+\norm{Tv_n}^2=s_1^2+\cdots+s_n^2$.
\item Prove that if $W=V$ and $T$ is a positive operator, then
  \[ \ip{Tv_1}{v_1}+\cdots+\ip{Tv_n}{v_n} = s_1+\cdots+s_n. \]
\end{parts}
\exnote{See the comment after Exercise 5 in Section 7A.}
\end{exercise}

\begin{exercise}\leavevmode
\begin{parts}
\item Give an example of a finite-dimensional vector space and an operator $T$
  on it such that the singular values of $T^2$ do not equal the squares of the
  singular values of $T$.
\item Suppose $T\in\Lin(V)$ is normal. Prove that the singular values of $T^2$
  equal the squares of the singular values of $T$.
\end{parts}
\end{exercise}

\begin{exercise}
Suppose $T_1,T_2\in\Lin(V)$. Prove that $T_1$ and $T_2$ have the same singular
values if and only if there exist unitary operators $S_1,S_2\in\Lin(V)$ such
that $T_1=S_1T_2S_2$.
\end{exercise}

\begin{exercise}
Suppose $T\in\Lin(V,W)$. Let $s_n$ denote the smallest singular value of $T$.
Prove that $s_n\norm{v}\le\norm{Tv}$ for every $v\in V$.
\end{exercise}

\begin{exercise}
Suppose $T\in\Lin(V)$ and $s_1\ge\cdots\ge s_n$ are the singular values of $T$.
Prove that if $\lambda$ is an eigenvalue of $T$, then
$s_1\ge\abs{\lambda}\ge s_n$.
\end{exercise}

\begin{exercise}
Suppose $T\in\Lin(V,W)$. Prove that $(T^*)^\dagger=(T^\dagger)^*$.
\exnote{Compare the result in this exercise to the analogous result for
invertible linear maps [see \ref{ax:7.5}(f)].}
\end{exercise}

\begin{exercise}
Suppose $T\in\Lin(V)$. Prove that $T$ is self-adjoint if and only if
$T^\dagger$ is self-adjoint.
\end{exercise}

\axquote{Matrices unfold\\ Singular values gleam like stars\\ Order in chaos
shines}{written by ChatGPT with input \emph{haiku about SVD}}
